\subsection{双代数与斜双代数}

定义3. 环A上的分次双代数（相应地，斜分次双代数）是一个集合E，它带有N型的分次A-代数结构和N型的分次A-余代数结构，且两者具有相同的底层分次A-模结构，并满足：

\begin{enumerate}
\def\labelenumi{(\arabic{enumi})}
\item
  A-代数 E 是结合的且有单位元的。
\item
  A-余代数 E 是余结合的且有余单位。
\item
  余积 \(c: \mathbf{E} \to \mathbf{E} \otimes_{\mathbf{A}} \mathbf{E}\) 是从分次代数 \(\mathbf{E}\) 到分次代数 \(\mathbf{E} \otimes_{\mathbf{A}} \mathbf{E}\) （相应地，分次代数 \(\mathbf{E}^{\mathfrak{g}} \otimes_{\mathbf{A}} \mathbf{E}\) ，参见§4，第7号）中的同态。
\item
  余单位 \(\gamma\) （E 的）是从分次代数 E 到代数 A（带平凡分次）中的分次代数同态，使得若 e 表示 A-代数 \(\mathbf{E}, \gamma(e) = 1\) 。
\end{enumerate}

如果 E 是一个分次双代数且其分次是平凡的，则 E 简称为双代数。一个分次双代数称为交换的（相应地，余交换的），如果其底代数是交换的（相应地，如果其底余代数是余交换的）；一个斜分次双代数称为反交换的（相应地，余反交换的），如果其底分次代数是反交换的（相应地，如果其底分次余代数是余反交换的）。

由定义3及第2节命题5可知，对于分次双代数或斜分次双代数E，

\[
\left\{ \begin{array}{c} c (e) = e \otimes e \\ c (x) \equiv x \otimes e + e \otimes x (\mathrm{mod.} \mathrm{E} _ {+} \otimes \mathrm{E} _ {+}) \quad \text { for } x \in \mathrm{E} _ {+} = \bigoplus_ {n \geqslant 1} \mathrm{E} _ {n}. \end{array} \right.\tag{22}
\]

若E和F是两个分次双代数（相应地，两个斜分次双代数），则一个映射 \(\phi: \mathrm{E} \to \mathrm{F}\) 称为分次双代数态射（相应地，斜分次双代数态射），如果：（1） \(\phi\) 是一个分次代数态射（因此将E的单位元映到F的单位元）；（2） \(\phi\) 是一个分次余代数态射，使得若 \(\gamma\) 和 \(\gamma'\) 分别是E和F的余单位，则 \(\gamma = \gamma' \circ \phi\) .

例子。（1）设 S 是一个具有单位元 u 的幺半群，于是幺半群 S 在 A 上的代数 \(E = A^{(S)}\) 具有单位元 \(e_{u}\) (§2，第6号)；另一方面已经看到，E 典范地具有一个带余单位 \(\gamma\) 的余结合、余交换 A-余代数结构，使得 \(\gamma(e_{s}) = 1\) （对所有 \(s \in S\) ）（第1号，例4及第2号，例4、例5和例6）。给出余积的公式 \(c(e_{s}) = e_{s} \otimes e_{s}\) 也立即表明 c 是一个代数同态。因此在 E 上定义了一个余交换双代数结构，带此结构的 E 被称为幺半群 S 在 A 上的双代数。

如果 T 是另一个带有单位元的幺半群， \(v, f: \mathrm{S} \to \mathrm{T}\) 为一个同态，使得 \(f(u) = v\) ，而 \(f_{(\mathbf{A})}: \mathrm{A}^{(\mathrm{S})} \to \mathrm{A}^{(\mathrm{T})}\) 为由 \(f(\S 2, \text{no. } 6)\) 导出的 A-代数同态，则可立即验证 \(f_{(\mathbf{A})}\) 是一个双代数同态。

\begin{enumerate}
\def\labelenumi{(\arabic{enumi})}
\setcounter{enumi}{1}
\item
  设 M 为一个 A-模。在 S(M) 上定义的分次 A-代数（§ 6，第 1 号）和分次 A-余代数（第 1 号，例 6）结构，在此集合上定义了一个交换余交换的分次双代数结构；因为我们已经看到（第 1 号，例 6），S(M) 上的余积是一个代数同态，并且由余单位 \(\gamma\) 的定义（第2号，例6）可知 \(\gamma(1) = 1\) ，且 \(\gamma\) 是从 \(E\) 到 \(A\) 中的代数同态。
\item
  设 M 为 A-模。我们如同例 2 中那样看到， \(\bigwedge\) (M) 上的分次 A-代数（§ 7，第 1 号）和分次 A-余代数（第 1 号，例 7）结构在此集合上定义一个反交换、余反交换的斜分次双代数结构。
\end{enumerate}

注：若 M 是一个 A-模，使得 \(M \otimes_{A} M \neq \{0\}\) ，则分次 A-代数（§ 5，第 1 号）和分次 A-余代数（第 1 号，例 5）在 T(M) 上的结构并不定义一个双代数结构，因为一般地

\[
c (x _ {1} x _ {2} y _ {1} y _ {2}) \neq c (x _ {1} x _ {2}) c (y _ {1} y _ {2})
\]

对 M 的四个元素 \(x_{1}, x_{2}, y_{1}, y_{2}\) 成立，如第1号公式 (3) 所示。

